% R582: EGA I, 0, 2.1.5--2.1.8.
% Traducción: OpenAI Codex, GPT-6.1 Sol, esfuerzo Ultra.
\oldpage[0\textsubscript{I}]{22}

\begin{env}[2.1.5]
\label{0.2.1.5-fr}
Sean $X$ un espacio irreducible y $f$ una aplicación continua de $X$ en un
espacio topológico $Y$. Entonces $f(X)$ es irreducible y, si $x$ es un punto
genérico de $X$, $f(x)$ es un punto genérico de $f(X)$ y, por tanto, también de
$\overline{f(X)}$. En particular, si además $Y$ es irreducible y tiene un solo
punto genérico $y$, para que $f(X)$ sea denso en todo el espacio, es necesario y suficiente
que $f(x)=y$.
\end{env}

\begin{env}[2.1.6]
\label{0.2.1.6-fr}
Todo subespacio irreducible de un espacio topológico $X$ está contenido en un
subespacio irreducible \emph{maximal}, que es necesariamente cerrado. Los
subespacios irreducibles maximales de $X$ se llaman \emph{componentes
irreducibles} de $X$. Si $Z_1$, $Z_2$ son dos componentes irreducibles
distintas del espacio $X$, $Z_1\cap Z_2$ es un conjunto cerrado \emph{de interior vacío}
en cada uno de los subespacios $Z_1$, $Z_2$; en particular, si una componente
irreducible de $X$ admite un punto genérico (2.1.2), tal punto no puede
pertenecer a ninguna otra componente irreducible. Si $X$ tiene un número
\emph{finito} de componentes irreducibles $Z_i$ ($1\leq i\leq n$) y si, para
cada $i$, se pone
\[
  U_i=Z_i\cap\complement\Bigl(\bigcup_{j\neq i}Z_j\Bigr),
\]
los $U_i$ son abiertos, irreducibles, disjuntos dos a dos y su
unión es densa en $X$.

Sea $U$ una parte abierta de un espacio topológico $X$. Si $Z$ es una parte
irreducible de $X$ que intersecta a $U$, $Z\cap U$ es abierto y denso en $Z$,
y, por tanto, irreducible; recíprocamente, para toda parte cerrada irreducible $Y$ de
$U$, la adherencia $\overline Y$ de $Y$ en $X$ es irreducible y
$\overline Y\cap U=Y$. Se concluye que hay una \emph{correspondencia
biyectiva} entre las componentes irreducibles de $U$ y las componentes
irreducibles de $X$ que intersectan a $U$.
\end{env}

\begin{env}[2.1.7]
\label{0.2.1.7-fr}
Si un espacio topológico $X$ es unión de un número \emph{finito} de
subespacios irreducibles cerrados $Y_i$, las componentes irreducibles de $X$
son los elementos maximales del conjunto de los $Y_i$, pues, si $Z$ es una parte
cerrada irreducible de $X$, $Z$ es unión de los $Z\cap Y_i$, de donde se deduce
inmediatamente que $Z$ debe estar contenido en uno de los $Y_i$. Sea $Y$ un subespacio
de un espacio topológico $X$ y supongamos que $Y$ tiene un número finito de
componentes irreducibles $Y_i$ ($1\leq i\leq n$); entonces las adherencias
$\overline{Y}_i$ en $X$ son las componentes irreducibles de $\overline Y$.
\end{env}

\begin{env}[2.1.8]
\label{0.2.1.8-fr}
Sea $Y$ un espacio irreducible que admite un solo punto genérico $y$. Sean
$X$ un espacio topológico y $f$ una aplicación continua de $X$ en $Y$.
Entonces, para toda componente irreducible $Z$ de $X$ que intersecta a $f^{-1}(y)$,
$f(Z)$ es denso en $Y$. La recíproca no es necesariamente cierta;
sin embargo, si $Z$ admite un punto genérico $z$ y si $f(Z)$ es denso en
$Y$, necesariamente se tiene $f(z)=y$ (2.1.5); además, $Z\cap f^{-1}(y)$ es
entonces la adherencia de $\{z\}$ en $f^{-1}(y)$ y, por tanto, es irreducible y,
como toda parte irreducible de $f^{-1}(y)$ que contiene a $z$ está
necesariamente contenida en $Z$ (2.1.6), $z$ es un punto genérico de
$Z\cap f^{-1}(y)$. Como toda componente irreducible de $f^{-1}(y)$ está
contenida en una componente irreducible de $X$, se ve que, si toda
componente irreducible $Z$ de $X$ que intersecta a $f^{-1}(y)$ admite un punto
genérico, hay una \emph{correspondencia biyectiva} entre el conjunto de
estas componentes y el conjunto de las componentes irreducibles
$Z\cap f^{-1}(Y)$ de $f^{-1}(y)$, siendo los puntos genéricos de $Z$
idénticos a los de $Z\cap f^{-1}(y)$.
\end{env}
